% Part: second-order-logic
% Chapter: syntax-and-semantics
% Section: introduction

\documentclass[../../../include/open-logic-section]{subfiles}

\begin{document}

\olfileid{sol}{syn}{int}

\olsection{پېژندنه}

په لومړۍ درجې منطق کښې د يوې ټاکلې ژبې غيرمنطقي سمبولونه، يعنې
!!{constant}s، !!{function}s او !!{predicate}s، له منطقي سمبولونو
سره يوځاے کوو څو د لومړۍ درجې !!{structure}s په باب خبرې بيانې کړو۔
دا کار د صدق د مفهوم له لارې کېږي: !!a{structure}~$\Struct{M}$، د
متغيرونو ګومارنه~$s$ او !!a{formula}~$!A$ سره تړي۔ $\Sat{M}{!A}[s]$
هله برقرار وي چې هغه څۀ رښتيا وي چې $!A$ يې وايي، کله چې د هغې
!!{constant}s، !!{function}s او !!{predicate}s د~$\Struct{M}$ له
مخې، او ازاد متغيرونه يې د~$s$ له مخې تفسير شي۔ د !!{identity}~$\eq$
تفسير د $\Sat{M}{!A}[s]$ په تعريف کښې دننه دے، او د~$\lforall$ او
$\lexists$ تفسير هم همداسې دے۔ لومړے تل د !!{structure} د
!!{domain}~$\Domain{M}$ پر سر د عينيت اړيکه ده؛ کميت ټاکونکي تل د
ټولې !!{domain} پر غړو ګرځي۔ خو مهمه دا ده چې کميت ټاکنه يوازې د
!!{domain} پر غړو روا ده؛ له کميت ټاکونکي وروسته يوازې د څيز
!!{variable}s راتلے شي۔

په دويمې درجې منطق کښې ژبه او د صدق تعريف دواړه غځېږي: ازاد او تړلي
د تابع او محمول متغيرونه، او پر هغو کميت ټاکنه، هم پکې راځي۔ دا
متغيرونه له !!{function}s او !!{predicate}s سره هماغه اړيکه لري چې
د څيز متغيرونه يې له !!{constant}s سره لري۔ د دويمې درجې منطق د
ترمونو او !!{formula}s په جوړولو کښې هماغسې ونډه لري، او کميت ټاکنه
پرې په ورته ډول کېږي۔ په \emph{معياري} معنٰی پوهنه کښې د دويمې درجې
کميت ټاکونکي د اړوند ډول پر ټولو ممکنو څيزونو ګرځي: د تابع متغيرونو
دپاره له $\Domain{M}$ څخه~$\Domain{M}$ ته پر $n$-ځايه تابعو، او د
محمول متغيرونو دپاره پر $n$-ځايه اړيکو۔ د بېلګې په توګه،
$\lforall[\Obj{v_0}][(\Obj{P^1_0}(\Obj{v_0}) \lor \lnot
  \Obj{P^1_0}(\Obj{v_0}))]$ د لومړۍ او دويمې درجې دواړو منطقونو
فارمولا ده؛ خو په دويمې درجې منطق کښې دا هم کتلے شو:
$\lforall[\Obj{V^1_0}][\lforall[\Obj{v_0}][(\Obj{V^1_0}(\Obj{v_0}) \lor
    \lnot \Obj{V^1_0}(\Obj{v_0}))]]$ او
$\lexists[\Obj{V^1_0}][\lforall[\Obj{v_0}][(\Obj{V^1_0}(\Obj{v_0}) \lor
    \lnot \Obj{V^1_0}(\Obj{v_0}))]]$۔ ځکه چې په دغو کښې ازاد
متغيرونه نشته، دا د دويمې درجې منطق !!{sentence}s دي۔ دلته
$\Obj{V^1_0}$ يو دويمې درجې $1$-ځايه محمول متغير دے۔ د
$\Obj{V^1_0}$ روا تفسيرونه هماغه دي چې د $\Obj{P^1_0}$ په شان
يو $1$-ځايه !!{predicate} ته يې ټاکلے شو، يعنې د~$\Domain{M}$
فرعي سټونه۔ نو پر هغو کميت ټاکنه دا وايي چې
$\lforall[\Obj{v_0}][(\Obj{V^1_0}(\Obj v_0) \lor \lnot
  \Obj{V^1_0}(v_0))]$ د~$\Domain{M}$ د يو فرعي سټ د
$\Obj{V^1_0}$ د ارزښت په توګه د ټاکلو په ټولو لارو کښې برقرار دے،
يا لږ تر لږه په يوه لاره کښې برقرار دے۔ ځکه چې هر سټ يو ورکړے څيز
يا لري يا يې نۀ لري، دواړه جملې په هر جوړښت کښې رښتيا دي۔
\end{document}
